problem:  
Let \( M^{4n} \) be a closed, simply connected Riemannian manifold with **positive sectional curvature**, endowed with an **effective isometric torus action** \( T^k \) of intermediate rank \( 2 < k < \lfloor n/2 \rfloor \).  
Assume that for every circle subgroup \( S^1 \subset T^k \), the fixed-point set \( M^{S^1} \) contains at least one totally geodesic submanifold \( N_{S^1} \subset M \) of **codimension exactly \(8\)**, and no circle-fixed-point component has smaller codimension.  

A manifold \(M\) is said to have **8‑periodic rational cohomology** if there exists \(x \in H^8(M;\mathbb{Q})\) such that multiplication by \(x\) induces isomorphisms  
\[
H^i(M;\mathbb{Q}) \xrightarrow{\cdot x} H^{i+8}(M;\mathbb{Q})
\]
for all \( i \le \dim M - 8.\)  

**Question:**  
Does the geometric and symmetry configuration above force \(M\) to have 8‑periodic rational cohomology?  
Equivalently, can there exist a positively curved, simply connected \(4n\)-manifold with this intermediate torus symmetry and only codimension‑8 circle-fixed submanifolds whose rational cohomology ring **fails** to be 8‑periodic?

Why is it a "good" problem:  
This problem isolates the precise threshold at which known curvature–symmetry theorems currently stop. Wilking’s connectedness lemma and Kennard’s periodicity theorem together guarantee 4‑periodicity when certain fixed-point or intersection submanifolds have codimension at most 6, but these arguments become inconclusive at codimension 8. The hypothesis that every \(S^1\)-fixed component has codimension exactly 8—with none smaller—precisely targets the **boundary between known periodicity and the unknown**.  

The question is thus not a generic restatement of the Hopf conjecture but a carefully sharpened inquiry into whether **new cohomological behaviors** can emerge under slightly weaker symmetry-induced constraints. Resolving it would require new insights linking **Riemannian geometry (positive curvature and totally geodesic fixed-point sets)**, **transformation group theory**, and **rational homotopy theory**.  

Because codimension‑8 symmetry-induced submanifolds are the first regime beyond the reach of existing tools, any advance here would clarify whether the rigidity phenomena characteristic of positive curvature persist further or whether qualitatively new topology can occur—making this a clean, concrete, and profoundly significant open problem.